\chapter{Conclusiones}
\label{cap:conclusiones}

En este capítulo se presentará una valoración global del trabajo realizado y de los resultados obtenidos. A diferencia de la discusión (Sección~\ref{sec:exp-discusion}), no se profundizará en la interpretación de los resultados experimentales, sino que se analizarán en relación con los objetivos establecidos en la Sección~\ref{sec:objetivos}. Asimismo, se propondrán posibles líneas de trabajo futuro que permitan ampliar y mejorar el sistema desarrollado.

\section{Conclusiones}
\label{sec:conc-conclusiones}

En esta sección se valorará el grado de cumplimiento de los objetivos específicos del trabajo, tomando como referencia la revisión del estado del arte (Capítulo~\ref{cap:estado-arte}), el análisis del conjunto de datos (Sección~\ref{sec:datos}), los métodos desarrollados (Sección~\ref{sec:metodos}) y los resultados experimentales (Capítulo~\ref{cap:experimentacion}). Para ello, se distinguirá entre el cumplimiento de los objetivos y la confirmación de la hipótesis planteada (Sección~\ref{sec:exp-numeracion}), dado que un resultado experimental desfavorable no implica necesariamente que un objetivo no se haya alcanzado. Asimismo, se diferenciarán las propiedades garantizadas por la asignación global de N3, como la unicidad de las etiquetas FDI, de aquellas que requieren una comprobación experimental, especialmente la identificación correcta de las posiciones correspondientes a dientes ausentes. Por último, se delimitará la aportación del trabajo al ámbito de la identificación forense, centrada en la segmentación y numeración coherente de los registros dentales tridimensionales como etapa previa a la comparación entre datos \textit{antemortem} y \textit{postmortem}, sin realizar directamente la identificación de personas.

\section{Trabajo futuro}
\label{sec:conc-trabajo-futuro}

En esta sección se propondrán posibles líneas de investigación que permitan abordar las limitaciones identificadas en la Sección~\ref{sec:exp-discusion}. En primer lugar, se planteará el desarrollo de un mapa de posiciones adaptable a la anatomía de cada arcada, con el fin de mejorar el etiquetado ante desplazamientos dentales producidos por la migración hacia posiciones vacantes. Asimismo, se considerará la ampliación de la evaluación a pacientes adultos con pérdidas dentales reales, dientes dañados y situaciones de etiquetado ambiguo, así como la extensión del sistema a la dentición temporal y mixta. Por último, se explorarán posibles aplicaciones en el ámbito forense, como la comparación automática entre registros \textit{antemortem} y \textit{postmortem} y el estudio del tamaño mínimo de fragmento dental necesario para realizar una identificación fiable, cuestión señalada por Forrest \cite{forrestForensicOdontologyDVI2019}.